\scanpage{70}\mspage{61}

il reste à voir que
\[
 |X|^\circ\cap|\partial Y|=\varnothing
\]
ce qui revient à
\[
 |X|^\circ\cap|Z|^{\deleted{\unsure{\circ}}}=\varnothing\qquad \forall Z\deleted{\unsure{\in}}\!<Y
\]
\deleted{\unsure{et si on}} \deleted{\unsure{(p}} et comme
\[
 |Z|=\bigcup_{Z'\leq Z}|Z'|^\circ
\]
par (iv), il revient au \unsure{$=$|même} de voir
\[
 |X|^\circ\cap|Z|^\circ=\varnothing\qquad\forall Z<Y
\]
ce qui va résulter de
\[
 X\,|\unsure{\circ|0}|\,Z\qquad\text{(via (ii))}
\]
or si $Z<Y$ on a $Z\,|\circ|\,Y$ ($\unsure{M_\theta\,4|M_\theta^{4}|\mathrm{Mag}\,4}$) et comme $X\mathrel{\overset{\circ}{\ll}}Z$ on en déduit $X\,|\circ|\,Z$ ($M_\theta\unsure{\,3|{}^{3}}$), on gagne !

Donc

\uline{Théorème} Les hom. de magasins$^{*}$
% slanted
\margin{cf p. 27 pour la \unsure{notion|notation}}
\[
 M\longrightarrow\mathrm{Figures}\,(P)\ \unsure{*|,}
\]
sont en corr. \unsure{1-1|1=1} avec les applications
\[
 \begin{aligned}
 M&\longrightarrow \unsure{\mathcal P|P}(P)\\
 X&\mathrel{\unsure{\text{---}|\to}}|X|
 \end{aligned}
\]
satisfaisant \unsure{:|à} les \deleted{applications} \inserted{conditions} (i)\,--\,(iv) de la page \deleted{\unsure{6}}59.

Une fonction ayant ces propriétés est
% archive sheet number: 70
